\section{Shared exponent generators and a strict improvement of the profile envelope}
\label{sec:sg-start}
\noindent\textbf{Status.} This continuation proves new ordinary mathematical
results using established logarithmic-form estimates. It gives unconditional
balance and small-prime estimates on actual content-one families, and a
conditional $abc$ class with a visibly retained large-prime hypothesis. It does
not prove or disprove standard $abc$. The elementary two-column reconstruction
has a separate Lean verification; the analytic theorems below are not thereby
kernel-checked.

Write $T=abc$, $R=\operatorname{rad}(T)$, $t=\log c$, $c=a+b$, where $a,b$ are
positive coprime integers. Retain the actual factorization from the preceding
exponent-profile supplement:
\[
 z=a+b\zeta=u(1+\zeta)^{e_0}\prod_{i=1}^r\pi_i^{e_i},\qquad
 M=N(z)=3^{e_0}\prod_{i=1}^r q_i^{e_i},
 \quad q_i=N(\pi_i)\ge7.
\]
The $q_i$ are distinct split primes, $e_0\in\{0,1\}$, and each split prime is
used with only one conjugate orientation. The prior identities are
\begin{equation}\label{eq:sg-baseline}
 \gcd(M,T)=1,\qquad \frac34c^2\le M\le c^2,\qquad
 z^3-\bar z^3=3\sqrt{-3}\,T.
\end{equation}
The case with no split factors is $(1,1,2)$ and satisfies $c\le R$ directly.
It is excluded from all logarithmic-form applications below.

\subsection{Allowing the same oriented prime in several generators}
\begin{definition}\label{def:sg-matrix}
A shared exponent representation has $m\ge1$ columns, integers $k_j\ge1$,
and nonnegative integers $f_{ij}$ such that
\[
 e_i=\sum_{j=1}^m k_j f_{ij},\qquad
 w_j=\prod_i\pi_i^{f_{ij}},\qquad Q_j=N(w_j),\qquad S=\sum_jk_j.
\]
Every column has a positive entry. Different columns may contain the same
oriented prime. The ramified factor stays outside the matrix.
\end{definition}

\begin{lemma}[Exact shared reconstruction]\label{lem:sg-reconstruction}
Every such representation satisfies
\[
 z=u(1+\zeta)^{e_0}\prod_jw_j^{k_j},\qquad
 \sum_j k_j\log Q_j=\log M-e_0\log3\le2t,\qquad Q_j\ge7.
\]
\end{lemma}
\begin{proof}
Expand the powers and commute the finite products. The total exponent of
$\pi_i$ is $\sum_j k_jf_{ij}=e_i$. Apply the multiplicative norm and take
logarithms. A nonempty column contains a prime of norm at least seven.
\end{proof}

\begin{theorem}[Shared-generator bounds at both kinds of places]
\label{thm:sg-two-place}
There is an effective absolute $C>2$ such that, for every representation,
\[
 \mathcal B=C^{m+1}\log(3+S)\prod_j\log Q_j
\]
satisfies, with $\lambda=(z/\bar z)^3$,
\begin{align}
 |\lambda-1|&\ge e^{-\mathcal B},&
 v_p(T)&\le p^2\mathcal B\quad(p\mid T),\label{eq:sg-local}\\
 \log\prod_{\substack{p\mid T\\p\le Y}}p^{v_p(T)}
 &\le\mathcal B Y^3\log Y\quad(Y\ge2),\label{eq:sg-small}\\
 c^3&\le8e^{\mathcal B}R^3\frac{E_3(T)}{C_3(T)}.\label{eq:sg-cubic}
\end{align}
Here $E_3(T)=\prod_{p\mid T}p^{(v_p(T)-3)_+}$ and
$C_3(T)=\prod_{p\mid T}p^{(3-v_p(T))_+}$.
\end{theorem}
\begin{proof}
Put $\eta_j=(w_j/\bar w_j)^3$. Both complex conjugates of $w_j$ have
absolute value $\sqrt{Q_j}$, so $h(w_j)=\tfrac12\log Q_j$ and the Weil
height inequalities imply $h^*(\eta_j)\le3\log Q_j$. No $\eta_j$ equals
one: otherwise the prime-ideal factorizations of $w_j^3$ and $\bar w_j^3$
would agree, although a selected split prime occurs only on the first side.
This argument needs no coprimality between different columns.

The unit contribution disappears and the ramified factor contributes
$(-1)^{e_0}$; hence $\lambda=(-1)^{e_0}\prod_j\eta_j^{k_j}\ne1$ by
\eqref{eq:sg-baseline}. Choose principal logarithms and an integer $v$ so that
\[
 \Lambda=\sum_jk_j\log\eta_j+(e_0-2v)\log(-1)
\]
has imaginary part in $[-\pi,\pi]$. It is nonzero and
$|e_0-2v|\le S+1$. Apply Bugeaud's Theorem 1.1, equation (1.2)
\cite{EPBugeaud2022}, with at most $m+1$ logarithms and field degree at most
two. Delete zero coefficients if needed and select a nonzero last coefficient.
Use $|e^{ix}-1|\ge2|x|/\pi$ on the displayed interval. Absorb the height
factors and fixed constants into one absolute $C$.

For $p\mid T$, the norm coprimality in \eqref{eq:sg-baseline} shows that
every $w_j,\bar w_j,z,\bar z$ is a unit at each place above $p$, including
when the $w_j$ overlap. With valuation normalized by $v_p(p)=1$,
\[
 v_p(\lambda-1)=v_p(T)+v_p(3\sqrt{-3})\ge v_p(T).
\]
Apply the first inequality of Bugeaud's Theorem 1.3 to
$\eta_1,\ldots,\eta_m,-1$ and coefficients $k_1,\ldots,k_m,e_0$.
It permits zero coefficients and does not require multiplicative independence.
Its $p^D$ factor is at most $p^2$. This proves \eqref{eq:sg-local} after
enlarging $C$ once; summation with
$\sum_{p\le Y}p^2\log p\le Y^3\log Y$ gives \eqref{eq:sg-small}.
Finally \eqref{eq:sg-baseline} gives
$T\ge c^3e^{-\mathcal B}/8$. Combine this with the exact identity
$TC_3(T)=R^3E_3(T)$ to prove \eqref{eq:sg-cubic}.
\end{proof}

Only the two explicitly identified established inequalities are used. In
particular the conditional refinements and open improvements later in the
non-Archimedean source are not substituted for its first inequality.

\subsection{Quotients, remainders, and a new uniform class}
Choose $g\ge1$ and nonnegative profiles $f_i,r_i$ with
$e_i=gf_i+r_i$ and some $f_i>0$. Set
\[
\begin{gathered}
 w=\prod_i\pi_i^{f_i},\qquad v=\prod_i\pi_i^{r_i},\qquad
 Q=N(w)\ge7,\qquad V=N(v),\\
 h=\max(1,\log V),\qquad \rho=\frac{h\log(4+g)}g.
\end{gathered}
\]
Then $z=u(1+\zeta)^{e_0}vw^g$. If $v$ is a unit, its column is omitted.
The parameter $g$ need not divide the original exponents.

\begin{corollary}\label{cor:sg-rho}
There is an effective absolute $A\ge1$ for which an admissible penalty satisfies
$\mathcal B\le A\rho t$. If $0<\rho\le1/64$ and $Y=\rho^{-1/6}$, then
\[
 \log\prod_{\substack{p\mid T\\p\le Y}}p^{v_p(T)}
 \le\frac A6\sqrt\rho\log(1/\rho)t.
\]
\end{corollary}
\begin{proof}
When $V>1$, use
$\mathcal B=C^3\log(4+g)\log V\log Q$ and $g\log Q\le2t$.
When $V=1$, the one-column bound is smaller than the same estimate with
$h=1$. Thus $A=2C^3$ suffices. Substitute
$Y^3=\rho^{-1/2}$ and $\log Y=\log(1/\rho)/6$ in
\eqref{eq:sg-small}; the threshold ensures $Y\ge2$.
\end{proof}

\begin{theorem}[Unconditional balance]\label{thm:sg-balance}
For every $0<\eta<1$ there is an effective $\rho_\eta>0$ such that every
triple with this representation and $\rho\le\rho_\eta$ satisfies
$\min(a,b)>c^{1-\eta}$. No hypothesis on its radical or large-prime tail is
required.
\end{theorem}
\begin{proof}
If $\min(a,b)\le c^{1-\eta}$, then $T\le c^{3-\eta}$ and
$\eta t\le\log8+\mathcal B$. Also $t\ge g\log7/2$ and $\rho\ge1/g$.
Consequently
$\eta\le(A+2\log8/\log7)\rho$, a contradiction after choosing a smaller
positive threshold.
\end{proof}

\begin{theorem}[An explicit conditional $abc$ class]\label{thm:sg-tail}
For every $B_0\ge0$ and $\varepsilon>0$ there is an effective
$\rho_{B_0,\varepsilon}>0$ such that
\[
 0<\rho\le\rho_{B_0,\varepsilon},\qquad
 W_{\rho^{-1/6}}(T):=\frac{E_3^{>\rho^{-1/6}}(T)}
 {C_3^{>\rho^{-1/6}}(T)}\le g^{B_0}
 \quad\Longrightarrow\quad c\le R^{1+\varepsilon}.
\]
The threshold is independent of the norm primes and their number.
\end{theorem}
\begin{proof}
Make the threshold at most $1/64$. The small-prime part of
$\log(E_3/C_3)$ is at most the full mass in Corollary~\ref{cor:sg-rho}, and
the remaining part is at most $B_0\log g$. Since $t\ge g\log7/2$ and
$\log g\le\log(4+g)$, the remainder in \eqref{eq:sg-cubic}, divided by
$t$, is at most
\[
 A\rho+\frac A6\sqrt\rho\log(1/\rho)
 +\frac{2(\log8+B_0)}{\log7}\rho.
\]
Choose the threshold to make this at most
$3\varepsilon/(1+\varepsilon)$. Taking logarithms in
\eqref{eq:sg-cubic} and rearranging gives
$t\le(1+\varepsilon)\log R$.
\end{proof}
The signed tail is an explicit unproved condition. The theorem asserts no
global bound on it and no infinitude of the class satisfying that condition.

\subsection{A strict separation from all disjoint partition envelopes}
\begin{theorem}\label{thm:sg-strict-separation}
There are height-unbounded primitive positive triples with split-exponent gcd
one such that every old disjoint partition penalty satisfies
\[
 \mathcal B_{\mathcal P}\ge\frac{\log4}{8}\log M,
\]
whereas a shared representation with two columns has
$\mathcal B_{\rm shared}=o(\log c)$. These triples satisfy the unconditional
balance and small-prime estimates above. Their large-prime tail remains open.
\end{theorem}
\begin{proof}
Take $r$ distinct split primes $q_i\in[X,2X]$, put $L=\log X\ge\log7$,
$g=\lceil r^4L^2\rceil$, and $e_i=g+i$. Choose one orientation $\pi_i$ of
each and form $z=\prod_i\pi_i^{e_i}$. A rational prime dividing both
coordinates would force a conjugate factor to occur, so $z$ is primitive.
A primitive zero-boundary pair is a unit, whereas $N(z)>1$. Unit rotation
therefore gives an actual positive triple. For $r\ge2$ the exponent gcd is one.
Writing $V_{\rm tot}=\log M$, we have
\begin{equation}\label{eq:sg-total}
 grL\le V_{\rm tot}\le4grL\le8r^5L^3.
\end{equation}
For $w=\prod_i\pi_i$ and $v=\prod_i\pi_i^i$,
$\log N(v)\le2r^2L$, and hence
\[
 \rho\le\frac{2\log(4+g)}{r^2L}\longrightarrow0
\]
uniformly in $L\ge\log7$. Indeed
$\log(4+g)\le\log6+4\log r+2\log L$; use the lower bound on $L$ and
boundedness of $(\log L)/L$. Corollary~\ref{cor:sg-rho} proves the new
sublinear estimate.

Consider any old partition into $m$ blocks. For a block of size $h_j$, write
$k_j$ for its exponent gcd and $L_j=\log Q_j$. Always $L_j\ge L$.
If $h_j\ge2$, its indices are distinct and congruent modulo $k_j$; therefore
\[
 k_j\le\frac{r-1}{h_j-1},\qquad
 L_j\ge\frac{g h_j(h_j-1)L}{r}.
\]
If $m\le r/2$, a largest block has $h\ge r/m\ge2$, giving
$L_j\ge grL/(2m^2)$. Bound the other columns by $L$ and use
\eqref{eq:sg-total} to obtain
\[
 \frac{\mathcal B_{\mathcal P}}{V_{\rm tot}}
 \ge\frac{C^{m+1}\log4\,L^{m-1}}{8m^2}\ge\frac{\log4}{8}.
\]
The last inequality uses $C>2$, $L>1$, and $2^{m+1}\ge m^2$. This elementary
inequality holds directly for $m=1,2,3$ and propagates for $m\ge3$ because
$(m+1)^2\le2m^2$.

If $m>r/2$, use $L_j\ge L$ in every column:
\[
 \frac{\mathcal B_{\mathcal P}}{V_{\rm tot}}
 \ge\frac{C^{m+1}\log4\,L^{m-3}}{8r^5}.
\]
For $r\ge8$, this is at least
\[
 \frac{\log4}{4(\log7)^3}\frac{(2\log7)^{r/2}}{r^5},
\]
which tends to infinity independently of $X$. Thus all partitions, including
those whose number of columns grows with $r$, satisfy the required lower bound
once $r$ is large.

Finally such prime lists exist with $r\to\infty$. As a separately identified
established input, the explicit fixed-modulus estimate of
\href{https://arxiv.org/abs/1802.00085v3}{Bennett, Martin, O'Bryant, and Rechnitzer}
gives $|\theta(x;3,1)-x/2|<x/(160\log x)$ for $x\ge8\cdot10^9$.
Subtract the estimates at $2X$ and $X$ to get
$\theta(2X;3,1)-\theta(X;3,1)\ge X/3$ eventually. The interval consequently
contains at least $X/(3\log(2X))$ such primes. This supplies $r\to\infty$;
\eqref{eq:sg-total} supplies unbounded height. This completes the proof.
\end{proof}

The lower bound is for the previous explicit envelope, not for the true angle.
The new representation therefore improves the available estimate on genuine
arithmetic inputs, without contradicting the earlier theorems. Its conditional
tail conclusion is not silently promoted to an unconditional $abc$ result.

\subsection{Dependencies, review, and precise formal scope}
The ordinary proof chain is actual oriented factorization, nonnegative matrix
reconstruction, the two stated established logarithmic-form bounds, and then
uniform balance and small-prime control. The conditional $abc$ arrow requires
the independent signed-tail hypothesis in Theorem~\ref{thm:sg-tail}.
Prime availability in Theorem~\ref{thm:sg-strict-separation} is an additional
established input used only in that family's infinitude argument.

The companion \texttt{OverlapColumns.lean} proves reconstruction and the norm
identity for a finite list of actual integer pairs with two shared exponent
columns, as well as Euclidean quotient--remainder splitting of exponents. Its
signature imposes no disjoint-support assumption. It does not formalize the
external analytic estimates, the splitting classification, all asymptotic
absorption, or standard $abc$.

The remaining arithmetic node for this new class is
$W_{\rho^{-1/6}}(T)\le g^{B_0}$ or a comparably sufficient correlated bound.
Prime-norm terminal profiles remain outside $\rho\to0$. Existing parent
routes remain open, and no claim of a closed Lean term proving
\texttt{ABCConjecture} is made.
